\documentclass[10pt,a4paper]{amsart}
\usepackage[margin=19mm]{geometry}
\usepackage{amsmath,amssymb,mathtools}
\usepackage[scr=rsfs,cal=euler]{mathalfa}
\usepackage{xfrac}
\usepackage[unicode,hidelinks,bookmarks=false]{hyperref}
\newcommand{\ie}{i.e.\ }
\newcommand{\cf}{cf.\ }
\newcommand{\resp}{respectively\ }
\newcommand{\Subsection}{\subsection}
\providecommand{\nobreakdash}{\nobreak\hbox{-}}
\setlength{\emergencystretch}{2em}
\multlinegap0pt
\usepackage{mathtools}
\usepackage{amssymb}
\usepackage{bm}
\usepackage[shortlabels]{enumitem}
\newtheorem{theorem}{Theorem}
\newtheorem{lemma}{Lemma}
\newcommand{\BS}{{\rm BS}}
\newcommand{\Om}{\Omega}
\renewcommand{\P}{{\mathbb P}}
\def\R{{\mathbb R}}
\newcommand{\T}{\mathbb{T}}
\newcommand{\dd}{\hspace{2pt}\mathrm{d}}
\newcommand{\del}{\Delta}
\newcommand{\eps}{\varepsilon}
\newcommand{\gam}{\gamma}
\newcommand{\grad}{\nabla}
\newcommand{\half}{\frac{1}{2}}
\newcommand{\loc}{{\rm loc}}
\newcommand{\mc}{\mathcal}
\title{GLOBAL WELL-POSEDNESS FOR THE 2D STOCHASTIC HYPOVISCOUS NAVIER-STOKES EQUATIONS}
\author{Daniel Goodair, Floris Roodenburg}
\date{Source: arXiv:2608.27054v1 (CC BY 4.0)}
\begin{document}
\maketitle

\noindent\emph{Source and licence.} GLOBAL WELL-POSEDNESS FOR THE 2D STOCHASTIC HYPOVISCOUS NAVIER-STOKES EQUATIONS. Version arXiv:2608.27054v1 (CC BY 4.0), distributed by arXiv under the Creative Commons Attribution 4.0 International licence, http://creativecommons.org/licenses/by/4.0/, as recorded in the arXiv licence metadata for this exact version and shown on the abstract page https://arxiv.org/abs/2608.27054v1. Full licence notice: \texttt{licenses/arxiv-2608.27054-CC-BY-4.0.txt}.\par\smallskip

\noindent\textit{Editorial scope.} Counted unit: the article's lemma lemma:vorticity\_to\_velocity (source0.tex lines 1093--1095). The article's complete original proof of it is delivered with it (source0.tex lines 1097--1107, one \texttt{proof} environment of 11 lines). No mathematics is written, repaired or re-derived: the statement and the proof are the article's own lines, in the article's own notation, and no retained line is substituted or re-flowed.  Retained uncounted as the source-local context the proof needs: lines 390--397; lines 923--930; lines 944--959; lines 1082--1082.  Nothing was omitted from any retained span, so the package carries no omission record.  No retained line contains a citation, so the wrapper carries no bibliography and no bibliography entry.  No retained line contains a figure environment, an image-inclusion command or drawing code, so no image is delivered and the package declares no asset dependency.  Editorial formatting only, in addition: an AMS \texttt{amsart} preamble that restates the article's own declarations and macros used by the retained lines, namely the environments \texttt{theorem}, \texttt{lemma} on separate counters, together with the standard packages those lines need.  The retained environments print in this document as Theorem 1, Lemma 1 in order of appearance, whereas the article prints the same items with the numbering of its own full document.  The complete native source of the article as distributed in the arXiv e-print for this exact version is bundled unchanged as \texttt{sources/208679/source0.tex}.
\begin{equation}\label{eq:abstract_eq}
       \left\{
        \begin{aligned}
            &\dd u + A u \dd t = F(u)\dd t + G(u)\dd W_{\ell^2}&\text{ for }&t\geq 0,\\
            &u|_{t=0}=u_0,
        \end{aligned}
        \right.
\end{equation}

\begin{equation}\label{eq: secondvorticity}
       \left\{
        \begin{aligned}
            &\dd \xi + (-\del)^\gam \xi \dd t  = -(\BS\xi \cdot \grad)\xi \dd t + \sum_{n \geq 1} \left(\alpha_n\xi + (\nabla \alpha_n) \times \BS\xi\right)\dd W^n_t&\text{ for }&t\geq 0,\\
            &\xi|_{t=0}=\xi_0,
        \end{aligned}
        \right.
\end{equation}

\begin{theorem}[Local well-posedness of the stochastic hypoviscous Navier--Stokes equations in vorticity form]\label{thm:localWPvorticity}
    Let $\gam\in (\frac 12,1]$, $s\in [\gam, 2\gam)$. Furthermore, assume that $p,q\in (2,\infty)$ satisfy
      \begin{equation*}
        \max\{|s-1|, 3\gam-s-1\}<\frac 2q < 4\gam-s-1\quad \text{ and }\quad \frac{2}{p}+\frac{2}{q\gam}\leq 4-\frac{s+1}{\gam}.
    \end{equation*}
    Moreover, set $\kappa:=-1 + \frac{p}{2}\big(4-\gam^{-1}(\frac{2}{q}+s+1)\big)$. Then for all $\xi_0 \in L^0_{\mc{F}_0}(\Om; \dot{B}_{q,p}^{\frac{2}{q}+1-2\gam}(\T^2;\R))$ there exists a unique maximal solution $(\xi,\Theta)$ to \eqref{eq: secondvorticity} satisfying $\Theta>0$ a.s. and
    \begin{align*}
        \xi&\in H^{\theta,p}_{\loc}([0,\Theta),w_{\kappa}; \dot{H}^{-s+2\gam(1-\theta),q}(\T^2; \R)) \,\,\text{a.s. for all }\theta\in [0,1/2),\\
        \xi&\in C([0, \Theta); \dot{B}_{q,p}^{\frac{2}{q}+1-2\gam}(\T^2; \R))\,\, \text{ a.s.}
    \end{align*}
    The solution $(\xi,\Theta)$ instantaneously regularises in time and space:
    \begin{equation*}
    \xi\in L^r_\loc((0,\Theta); \dot{H}^{\gam,\zeta}(\T^2;\R))\cap C^{\theta -\eps}_\loc((0,\Theta); \dot{H}^{-\gam+2\gam(1-\theta),\zeta}(\T^2;\R)) \quad \text{a.s.},
\end{equation*}
for $\theta \in (0,\half)$, $\eps\in (0,\theta)$,  $r,\zeta\in (2,\infty)$.
\end{theorem}

\begin{equation} \label{global noise}g_n(x,u):= \Pi[\alpha_n(x)u]\quad \text{ with }(\alpha_n)_{n\geq 1}\in \ell^2(W^{1,\infty}(\T^2)).\end{equation}

\begin{lemma} \label{lemma:vorticity_to_velocity}
    Let $(\xi,\Theta)$ be the unique maximal solution to \eqref{eq: secondvorticity}, specified in Theorem \ref{thm:localWPvorticity}. Then $(\BS\xi, \Theta)$ is a local solution to \eqref{eq:abstract_eq} for the choice of $g_n$ given by \eqref{global noise}. Furthermore let $(u,\sigma)$ denote the unique maximal solution to \eqref{eq:abstract_eq} with the initial condition $u_0 \coloneqq \BS\xi_0$. Then $\Theta \leq \sigma$ a.s. and $u = \BS\xi$ on $[0,\Theta)$. 
\end{lemma}

\begin{proof}
    As $\BS\xi$ has only better regularity than $\xi$, we just need to verify that the identity is satisfied. We recall that for a zero-mean function $f$, $\BS f$ is the unique divergence-free and zero-mean vector field such that $\nabla \times (\BS f) = f$. Using that the equation \eqref{eq: secondvorticity} was obtained by taking the curl of $u$, we have already seen that
    \begin{align*}
        \nabla \times \left[(-\del)^\gam \BS\xi \right] &= (-\del)^\gam \xi,\\
        \nabla \times \left[\P \left((\BS\xi \cdot \grad)\BS\xi \right)\right] &= (\BS\xi \cdot \grad)\xi,\\
        \nabla \times \left[ \P \Pi \left(\alpha_n \BS\xi\right)\right] &= \alpha_n\xi + (\nabla \alpha_n) \times \BS\xi.
    \end{align*}
   Thus applying $\BS$ term by term to the identity satisfied by $\xi$, we see that on any localising sequence for $\Theta$, the identity
   $$\dd \BS\xi + [(-\del)^\gam \BS\xi+\P\left((\BS\xi\cdot\grad) \BS\xi\right)]\dd t  =\sum_{n\geq 1}\P \Pi(\alpha_n\BS\xi)\dd W^n_t$$
holds. Hence, $(\BS\xi, \Theta)$ is a local solution to \eqref{eq:abstract_eq}. The second part of the lemma now simply follows from the definition of the maximal time and uniqueness; maximality ensures that $\Theta \leq \sigma$ a.s., whilst uniqueness implies that $\BS\xi = u$ on $[0,\Theta)$.
\end{proof}

\end{document}
